% Memoria del trabajo de criptografía post-cuántica (CSD, MUCC - UPV)
% Requisitos del enunciado: máx. 5 páginas, fuente 12 pt, márgenes <= 2,54 cm, interlineado simple.
% Secciones obligatorias y peso: presentación 15 %, algoritmos 30 %, tareas 40 %, conclusiones 15 %.
% Compilar:  latexmk -pdf memoria.tex      (o:  python build.py)
\documentclass[12pt,a4paper]{article}

\usepackage[utf8]{inputenc}
\usepackage[T1]{fontenc}
\usepackage[spanish,es-tabla]{babel}
\usepackage[margin=2.54cm]{geometry}
\usepackage{amsmath,amssymb}
\usepackage{booktabs}
\usepackage{graphicx}
\usepackage{xcolor}
\usepackage[hidelinks]{hyperref}

\setlength{\parindent}{0pt}
\setlength{\parskip}{4pt}

% Marcador de pendientes (en rojo). Buscar "TODO" antes de entregar.
\newcommand{\todo}[1]{\textcolor{red}{[TODO: #1]}}

\title{\textbf{Criptoanálisis de criptografía post-cuántica:\\ GGH y LWE}}
\author{\todo{Nombre 1, Nombre 2, Nombre 3, Nombre 4} \\[2pt]
        Criptología y Seguridad de los Datos \\
        Máster en Ciberseguridad y Ciberinteligencia, ETSInf-UPV}
\date{\today}

\begin{document}
\maketitle

%% ---------------------------------------------------------------------------
\section{Introducción}
% Peso 15 % junto con la presentación general.
Objetivo del trabajo y problemas abordados: \emph{Shortest Vector Problem} (SVP),
\emph{Closest Vector Problem} (CVP), cifrado GGH, \emph{Learning With Errors} (LWE) y Ring-LWE.
\todo{Tareas elegidas (tier y número) y criterio de elección. Lenguaje y entorno: Python 3 + numpy.}

%% ---------------------------------------------------------------------------
\section{Algoritmos}
% Peso 30 %. Por cada algoritmo: qué hace, POR QUÉ resuelve el problema, coste, hipótesis.
\subsection{Reducción de bases y aproximaciones de Babai}
\todo{LLL: idea, condiciones de reducción, por qué sirve contra GGH. Redondeo de Babai y plano más cercano.}

\subsection{SVP y CVP exactos}
\todo{Fuerza bruta (justificar cuándo es viable) y enumeración. Justificación de la cota de búsqueda.}

\subsection{LWE y Ring-LWE}
\todo{Estrategia para recuperar el secreto, descifrado de un bit (umbral), caso Ring-LWE en $\mathbb{Z}_q[x]/(x^n+1)$.}

\subsection{Aportaciones propias}
\todo{Si las hay: hipótesis en la que se basan y por qué se espera que funcionen.}

%% ---------------------------------------------------------------------------
\section{Resolución de tareas}
% Peso 40 %. Traza breve con resultados intermedios + evolución del tiempo con la talla.
\begin{table}[ht]
\centering
\begin{tabular}{llllr}
\toprule
Tarea & Tier & Método & Resultado & Tiempo (s) \\
\midrule
\todo{Tier0T01\_Latt} & 0 & \todo{} & \todo{} & \todo{} \\
\bottomrule
\end{tabular}
\caption{Resumen de tareas resueltas.}
\label{tab:tareas}
\end{table}

\todo{Traza de ejecución de al menos una tarea por algoritmo (resultados intermedios).}

% Gráfica de tiempo frente a talla: se incluye si existe results/bench.pdf
\IfFileExists{../results/bench.pdf}{%
  \begin{figure}[h]\centering
  \includegraphics[width=0.7\linewidth]{../results/bench.pdf}
  \caption{Tiempo de resolución en función de la talla del problema.}
  \end{figure}}{\todo{Gráfica tiempo vs.\ talla (scripts/bench.py $\to$ results/bench.pdf).}}

%% ---------------------------------------------------------------------------
\section{Conclusiones}
% Peso 15 %.
\todo{Qué se ha aprendido sobre la dificultad de los problemas, cómo escala el coste, limitaciones.}

%% ---------------------------------------------------------------------------
\begin{thebibliography}{9}
\bibitem{ggh} O.~Goldreich, S.~Goldwasser, S.~Halevi.
  \emph{Public-key cryptosystems from lattice reduction problems}. CRYPTO '97, pp.\ 112--131, 1997.
\bibitem{regev} O.~Regev.
  \emph{On lattices, learning with errors, random linear codes, and cryptography}.
  Journal of the ACM 56(6):34, 2009.
\bibitem{lll} A.~K.~Lenstra, H.~W.~Lenstra, L.~Lovász.
  \emph{Factoring polynomials with rational coefficients}. Math.\ Annalen 261, 1982.
\todo{Resto de recursos bibliográficos y tecnológicos utilizados.}
\end{thebibliography}

\end{document}
